\documentclass[11pt, a4paper]{article}

\usepackage[utf8]{inputenc}
\usepackage[T1]{fontenc}
\usepackage{amsmath, amssymb, amsthm}
\usepackage{graphicx}
\usepackage{hyperref}
\usepackage{booktabs}
\usepackage{geometry}
\usepackage{lmodern}

\geometry{margin=1in}

\hypersetup{
  colorlinks=true,
  linkcolor=blue,
  citecolor=blue,
  urlcolor=blue
}

\newtheorem{theorem}{Theorem}
\newtheorem{lemma}[theorem]{Lemma}

\title{Exploring the Foundations of Mathematical Analysis}
\author{cc-latex Demo Project}
\date{\today}

\begin{document}

\maketitle

\begin{abstract}
This document demonstrates the capabilities of the cc-latex IDE, including
multi-file projects, mathematical typesetting, tables, cross-references,
and bibliography management. Edit any file in the tree on the left, or ask
the AI assistant below for help writing your next paper.
\end{abstract}

\tableofcontents

\section{Introduction}
\label{sec:introduction}

\LaTeX{} is the standard document preparation system for academic and scientific
publishing. Originally created by Leslie Lamport as a set of macros on top of
Donald Knuth's \TeX{} typesetting engine, it has become indispensable for
producing high-quality technical documents.

\subsection{Why \LaTeX{}?}

There are several compelling reasons to use \LaTeX{} over conventional
word processors:

\begin{itemize}
  \item \textbf{Mathematical typesetting} --- unmatched precision and beauty
    for equations, from simple inline expressions like $E = mc^2$ to complex
    multi-line derivations.
  \item \textbf{Consistent formatting} --- separation of content and
    presentation ensures uniform styling throughout the document.
  \item \textbf{Cross-referencing} --- automatic numbering of sections,
    equations, figures, and tables with reliable cross-references.
  \item \textbf{Bibliography management} --- seamless integration with
    Bib\TeX{} for citation handling.
  \item \textbf{Version control} --- plain text source files work naturally
    with Git and other version control systems.
\end{itemize}

\subsection{About This Project}

This file (\texttt{introduction.tex}) is included via the \verb|\input{}|
command in \texttt{main.tex}, demonstrating multi-file project support.
You can split large documents across multiple files for better organization.

Try editing this section and recompiling to see your changes reflected in the
PDF preview on the right.


\section{Mathematical Typesetting}

One of \LaTeX's greatest strengths is its ability to render complex mathematics
with precision and clarity.

\subsection{Equations}

The Euler--Lagrange equation governs the calculus of variations:
\begin{equation}
  \frac{\partial \mathcal{L}}{\partial q} - \frac{d}{dt}\frac{\partial \mathcal{L}}{\partial \dot{q}} = 0
  \label{eq:euler-lagrange}
\end{equation}

A particularly elegant identity connects five fundamental constants:
\begin{equation}
  e^{i\pi} + 1 = 0
  \label{eq:euler}
\end{equation}

\subsection{Theorems}

\begin{theorem}[Fundamental Theorem of Calculus]
  Let $f$ be a continuous function on $[a, b]$ and let $F$ be an antiderivative
  of $f$. Then
  \[
    \int_a^b f(x) \, dx = F(b) - F(a).
  \]
\end{theorem}

\begin{lemma}
  For any $n \times n$ matrix $A$,
  \[
    \det(A^T) = \det(A).
  \]
\end{lemma}

\subsection{Aligned Equations}

Maxwell's equations in differential form:
\begin{align}
  \nabla \cdot \mathbf{E}  &= \frac{\rho}{\varepsilon_0}          \label{eq:gauss} \\
  \nabla \cdot \mathbf{B}  &= 0                                    \label{eq:gauss-mag} \\
  \nabla \times \mathbf{E} &= -\frac{\partial \mathbf{B}}{\partial t} \label{eq:faraday} \\
  \nabla \times \mathbf{B} &= \mu_0 \mathbf{J} + \mu_0 \varepsilon_0 \frac{\partial \mathbf{E}}{\partial t} \label{eq:ampere}
\end{align}

\section{Tables and Figures}

Table~\ref{tab:complexity} summarizes the time complexity of common sorting algorithms.

\begin{table}[ht]
  \centering
  \caption{Comparison of sorting algorithm complexities}
  \label{tab:complexity}
  \begin{tabular}{@{} l c c c @{}}
    \toprule
    Algorithm      & Best Case        & Average Case     & Worst Case       \\
    \midrule
    Merge Sort     & $O(n \log n)$    & $O(n \log n)$    & $O(n \log n)$    \\
    Quick Sort     & $O(n \log n)$    & $O(n \log n)$    & $O(n^2)$         \\
    Heap Sort      & $O(n \log n)$    & $O(n \log n)$    & $O(n \log n)$    \\
    Insertion Sort & $O(n)$           & $O(n^2)$         & $O(n^2)$         \\
    Bubble Sort    & $O(n)$           & $O(n^2)$         & $O(n^2)$         \\
    \bottomrule
  \end{tabular}
\end{table}

\section{Cross-References}

As shown in Equation~\ref{eq:euler-lagrange}, variational principles are
foundational in physics. Gauss's law (Equation~\ref{eq:gauss}) is the first
of Maxwell's equations. The sorting complexities in Table~\ref{tab:complexity}
highlight the efficiency of divide-and-conquer approaches.

\section{Conclusion}

This demo project showcases multi-file \LaTeX{} editing with cc-latex.
Try editing this file or \texttt{introduction.tex}, then compile with
\textbf{Ctrl+Shift+B} to see the result. Use the AI chat below for help
writing, formatting, or debugging your documents.

For more on \LaTeX{} best practices, see Knuth~\cite{knuth1984} and
Lamport~\cite{lamport1994}. Online resources like the Overleaf
documentation~\cite{overleaf2024} are also excellent references.

\bibliographystyle{plain}
\bibliography{references}

\end{document}
